\documentclass[journal, onecolumn]{IEEEtran}

\usepackage{cite} \usepackage{amsmath,amssymb,amsfonts} \usepackage{graphicx} \usepackage{textcomp} \usepackage{xcolor} \usepackage{booktabs} \usepackage{multirow} \usepackage{url} \usepackage{hyperref} \usepackage{array} \usepackage{tabularx} \usepackage{listings} \usepackage{microtype}

\widowpenalty=10000 \clubpenalty=10000

\hypersetup{ colorlinks=true, linkcolor=black, citecolor=blue, urlcolor=blue }

\lstset{
  basicstyle=\ttfamily\small,
  breaklines=true,
  frame=single,
  language=Python,
  numbers=left,
  numberstyle=\tiny\color{gray},
}

\begin{document}

\title{Optimizing a Python \& PostgreSQL Data Pipeline\\for Large-Scale JSONB Expansion over SSH}

\author{\IEEEauthorblockN{Pablo Alessandro Santos Hugen} \IEEEauthorblockA{Federal University of Rio Grande do Sul \\ Institute Of Informatics \\ Porto Alegre, RS, Brazil \\ pablo.hugen@inf.ufrgs.br}}

\maketitle

\begin{abstract}
TODO
\end{abstract}

\begin{IEEEkeywords}
PostgreSQL, data pipeline, JSONB, performance optimization, Python, SSH, database expansion, checkpoint recovery.
\end{IEEEkeywords}

\section{Introduction}
\label{sec:intro}

This work addresses the Python \& PostgreSQL Data Pipeline Optimization Challenge~\cite{grimm2026}.

TODO

\section{Architecture and Background}
\label{sec:architecture}

TODO

\section{Methodology and Results}
\label{sec:methodology}

TODO

\section{Conclusion}
\label{sec:conclusion}

TODO

\bibliographystyle{IEEEtran}
\bibliography{references}

\end{document}
